\section{Factor Demand \& Marshall's Laws}
Factor demand ties consumer theory to production: the same substitution-vs-scale decomposition that drives Slutsky drives Marshall's Laws. Every consumer result has a producer dual (see Duality Table in \S Math Toolkit).

\subsection{Derived Demand}
Demand for an input is \textit{derived} from demand for the output it produces. A wage increase $w \uparrow$ affects $L^*$ through two channels:

\textbf{Substitution effect} (hold output fixed, reoptimize inputs): $\partial L^c/\partial w < 0$.

\textbf{Scale effect} (higher cost $\implies$ higher price $\implies$ lower output): reinforces the substitution effect for normal inputs.

\subsection{Two-Input Model}
Firm produces $q = f(K,L)$, sells at price $p$. Cost minimization:
$\frac{f_L}{f_K} = \frac{w}{r}$.

\textbf{Elasticity of substitution:}
$\sigma = \frac{d\ln(K/L)}{d\ln(w/r)}\big|_{q\text{ fixed}}$

$\sigma = 0$: Leontief (no substitution). $\sigma = 1$: Cobb-Douglas. $\sigma = \infty$: linear technology.

\subsection{Marshall's Laws of Derived Demand}
Demand for a factor is \textbf{more elastic} when:
\begin{enumerate}[nosep]
    \item \textbf{Elasticity of substitution} $\sigma$ is high
    \item \textbf{Demand for output} $|\epsilon_D|$ is high
    \item Factor's \textbf{cost share} $s$ is large (caveat: when $\sigma < |\epsilon_D|$, larger share makes demand \textit{less} elastic)
    \item \textbf{Supply of other factors} $\epsilon_S^{other}$ is high
\end{enumerate}

\textbf{Two-factor formula:}
$$\eta_L = \frac{(1-s_L)\sigma\epsilon_S^K + s_L|\epsilon_D|\sigma + (1-s_L)|\epsilon_D|\epsilon_S^K}{(1-s_L)\epsilon_S^K + s_L\sigma + (1-s_L)|\epsilon_D|}$$

With $\epsilon_S^K = \infty$ (perfectly elastic capital supply):
$$\eta_L = (1-s_L)\sigma + s_L|\epsilon_D|$$
A weighted average: substitution effect (weight $1-s_L$) and scale effect (weight $s_L$).

\textbf{Cross-price elasticity:}
$\beta_{ij} = s_j\epsilon_D + s_j\sigma_{ij}$

\textit{Application (education production):} Inputs: student time ($sK$), purchased materials ($D$). Remote learning lowers productivity of $sK$ $\implies$ MC rises $\implies$ scale effect reduces all inputs. Substitution: shift toward purchased inputs (online tools). Decompose via Marshall's Laws: scale dominates when $\sigma$ is small, substitution dominates when $\sigma$ is large.

\textit{Application (min wage):} $w$ forced above $w^*$: substitution replaces workers with capital; scale reduces output. Total effect depends on $\sigma$ and $|\epsilon_D|$ via Marshall's formula.

\subsection{Minimum Wage}
\textbf{Competitive labor market:} floor $\bar w > w^*$ $\implies$ employment falls along demand ($\Delta L \approx \eta_L \cdot \% \Delta w$); unemployment $=$ excess supply, which \textit{exceeds} job loss (higher wage draws entrants into search).

\textbf{Non-wage margin adjustment} (the Chicago answer): firms recapture the mandated wage on unregulated margins---cut training, benefits, schedule flexibility, raise effort demands. Total compensation moves toward the market-clearing level $\implies$ measured \textit{employment} effects understate the distortion, and covered workers gain less than the wage bump suggests.

\textbf{Monopsony:} employer faces upward-sloping labor supply, $MC_L > w$. A floor set in $(w^m, MC_L^m)$ acts like the price ceiling on a monopolist: \textit{raises} employment. Empirical relevance where mobility is low (company towns, non-competes, nurses).

\textbf{Incidence \& selection:} job losses concentrate on the least skilled within the covered group (employers ration jobs by quality); some rents pass to consumers or landlords in low-wage regions.

\textit{TFU trap:} ``small measured disemployment $\implies$ minimum wage has no cost'' --- U: margins other than headcount (hours, training, amenities) absorb the shock; also monopsony gives the same observation with the opposite welfare sign. State which model and which margins.

\subsection{Isoprofit Curves}
For a firm choosing $(w, L)$ along an isoprofit: $\pi = pf(K,L) - wL - rK$. The isoprofit curve in $(L,w)$ space is:
$$w = \frac{pf(K,L) - rK - \bar{\pi}}{L}$$
Higher isoprofits are closer to the origin. The firm's labor demand curve traces out the tangencies of isoprofit curves with the wage line.
